% from ega4-11-fr.tex
La démonstration de (ii) se ramène de la même façon au cas où $S$ est
noethérien,
\oldpage[IV-3]{134}
en utilisant (11.2.9)~; la conclusion est alors évidente,
$\operatorname{gr}_{\mathfrak J}^\bullet(B)$ étant dans ce cas une
$(B/\mathfrak J)$-algèbre de type fini et
$\operatorname{gr}_{\mathfrak J}^\bullet(M)$ un
$\operatorname{gr}_{\mathfrak J}^\bullet(B)$-module de type fini.

Pour prouver (iii), on se ramène comme dans (i) au cas où $S$ et $X$ sont
affines~; avec les mêmes notations, on peut supposer, en vertu de (11.2.9) que
$\operatorname{gr}_{\mathfrak J_\lambda}^\bullet(B_\lambda)$ et
$\operatorname{gr}_{\mathfrak J_\lambda}^\bullet(M_\lambda)$ sont des
$A_\lambda$-modules plats et que l'on a
$\operatorname{gr}_{\mathfrak J}^\bullet(B)
=\operatorname{gr}_{\mathfrak J_\lambda}^\bullet(B_\lambda)\otimes_{A_\lambda}A$
et
$\operatorname{gr}_{\mathfrak J}^\bullet(M)
=\operatorname{gr}_{\mathfrak J_\lambda}^\bullet(M_\lambda)\otimes_{A_\lambda}A$.
Par suite $\operatorname{gr}_{\mathfrak J}^\bullet(B)$ est une
$(B/\mathfrak J)$-algèbre de présentation finie et
$\operatorname{gr}_{\mathfrak J}^\bullet(M)$ un
$\operatorname{gr}_{\mathfrak J}^\bullet(B)$-module de présentation finie. Si $W^m$
est l'ensemble des $\mathfrak p\in\operatorname{Spec}(B/\mathfrak J)$ tels que
$\operatorname{gr}_{\mathfrak J}^m(M)_\mathfrak p$ soit un
$(B/\mathfrak J)$-module plat, $W^m$ est ouvert dans
$\operatorname{Spec}(B/\mathfrak J)$ (11.3.1) et l'on a
$W=\bigcap_mW^m$, donc $W$ est stable par générisation. L'assertion (iii)
résulte alors de (11.3.2) appliqué en y prenant
$X=\operatorname{Spec}(\operatorname{gr}_{\mathfrak J}^\bullet(B))$,
$Y=Z=\operatorname{Spec}(B/\mathfrak J)$ et
$\mathcal F=(\operatorname{gr}_{\mathfrak J}^\bullet(M))^{\sim}$.

Enfin, (iv) découle aussitôt de (11.2.9.2).

Généralisant la définition de (6.10.1), on dit que pour un préschéma $X$, un
sous-préschéma fermé $Y$ de $X$ défini par un Idéal quasi-cohérent $\mathcal J$
de $\mathcal O_X$, et un $\mathcal O_X$-Module quasi-cohérent $\mathcal F$,
$\mathcal F$ est \emph{normalement plat le long de $Y$ en un point $y$} si
$(\operatorname{gr}_{\mathcal J}^\bullet(\mathcal F))_y$ est un
$\mathcal O_{Y,y}$-module \emph{plat}. On dit que $F$ est
\emph{normalement plat le long de $Y$} s'il l'est en tous les points de $Y$.
\end{proof}

\begin{corollary}[11.3.5]
\label{IV.11.3.5-fr}
Sous les hypothèses générales de \textup{(11.3.4)}, on suppose que
$\operatorname{gr}_{\mathcal J}^\bullet(\mathcal O_X)$ et
$\operatorname{gr}_{\mathcal J}^\bullet(\mathcal F)$ soient $f$-plats. Alors
l'ensemble $W$ des $y\in Y$ tels que $\mathcal F$ soit normalement plat le long
de $Y$ au point $y$ est ouvert dans $Y$, et (avec les notations de
\textup{(11.3.4, (iv)))} $\mathcal F'$ est normalement plat le long de $Y'$ en tous les
points de $W'=p^{-1}(W)$.
\end{corollary}

\begin{proposition}[11.3.6]
\label{IV.11.3.6-fr}
Les notations étant celles de \textup{(8.5.1)} et \textup{(8.8.1)}, on suppose $S_\alpha$
quasi-compact, $X_\alpha$ de présentation finie sur $S_\alpha$, $Y_\alpha$ un
sous-préschéma fermé de $X_\alpha$ défini par un Idéal quasi-cohérent
$\mathcal J_\alpha$ de type fini de $\mathcal O_{X_\alpha}$ tel que
$\operatorname{gr}_{\mathcal J_\alpha}^\bullet(\mathcal O_{X_\alpha})$ soit plat sur
$S_\alpha$~; enfin on suppose que $\mathcal F_\alpha$ est un
$\mathcal O_{X_\alpha}$-Module de présentation finie. Pour que $\mathcal F$
soit normalement plat le long de $Y$, il faut et il suffit qu'il existe
$\lambda$ tel que $\mathcal F_\lambda$ soit normalement plat le long de
$Y_\lambda$.
\end{proposition}

\begin{proof}
Notons que $Y$ (resp. $Y_\lambda$) est le sous-préschéma fermé de $X$ (resp.
$X_\lambda$) défini par $\mathcal J=\mathcal J_\alpha\mathcal O_X$ (resp.
$\mathcal J_\lambda=\mathcal J_\alpha\mathcal O_{X_\lambda}$)~; en vertu de
l'hypothèse et de (11.2.9.2), on a
$\operatorname{gr}_{\mathcal J}^\bullet(\mathcal O_X)
=\operatorname{gr}_{\mathcal J_\alpha}^\bullet(\mathcal O_{X_\alpha})
 \otimes_{\mathcal O_{S_\alpha}}\mathcal O_S$
et
$\operatorname{gr}_{\mathcal J_\lambda}^\bullet(\mathcal O_{X_\lambda})
=\operatorname{gr}_{\mathcal J_\alpha}^\bullet(\mathcal O_{X_\alpha})
 \otimes_{\mathcal O_{S_\alpha}}\mathcal O_{S_\lambda}$
pour $\lambda\geq\alpha$,
et $\operatorname{gr}_{\mathcal J}^\bullet(\mathcal O_X)$ (resp.
$\operatorname{gr}_{\mathcal J_\lambda}^\bullet(\mathcal O_{X_\lambda})$) est plat
sur $S$ (resp. $S_\lambda$), ce qui entraîne que $Y$ est plat sur $S$ (resp.
$Y_\lambda$ plat sur $S_\lambda$). Si $\mathcal F_\lambda$ est normalement plat
le long de $Y_\lambda$, $\operatorname{gr}_{\mathcal J_\lambda}^\bullet
(\mathcal F_\lambda)|Y_\lambda$ est plat sur $Y_\lambda$, donc aussi sur
$S_\lambda$, et comme
$(\operatorname{gr}_{\mathcal J_\lambda}^\bullet(\mathcal F_\lambda))_{x_\lambda}=0$
pour $x_\lambda\notin Y_\lambda$,
$\operatorname{gr}_{\mathcal J_\lambda}^\bullet(\mathcal F_\lambda)$ est plat sur
$S_\lambda$. On en conclut (11.2.9.2) que
$\operatorname{gr}_{\mathcal J}^\bullet(\mathcal F)
=\operatorname{gr}_{\mathcal J_\lambda}^\bullet(\mathcal F_\lambda)
 \otimes_{\mathcal O_{S_\lambda}}\mathcal O_S$,
donc $\mathcal F$ est normalement plat le long de $Y$. Pour prouver la
réciproque, on peut supposer que $S_\alpha$ et $X_\alpha$ sont affines et
adopter les notations de (11.2.9)~; puisque
$\operatorname{gr}_{\mathfrak J}^\bullet(B)$ est un $A$-module plat, il en est de même
de $\operatorname{gr}_{\mathfrak J}^\bullet(M)$, donc, en vertu de (11.2.9), il
existe $\lambda\geq\alpha$ tel que
$\operatorname{gr}_{\mathfrak J_\lambda}^\bullet(M_\lambda)$ soit un
$A_\lambda$-module plat, d'où
$\operatorname{gr}_{\mathfrak J}^\bullet(M)
=\operatorname{gr}_{\mathfrak J_\lambda}^\bullet(M_\lambda)\otimes_{A_\lambda}A$.
En outre (11.3.4, (ii)),
$\operatorname{gr}_{\mathfrak J_\lambda}^\bullet(M_\lambda)$ est un
$\operatorname{gr}_{\mathfrak J_\lambda}^\bullet(B_\lambda)$-module de présentation
finie et $\operatorname{gr}_{\mathfrak J_\lambda}^\bullet(B_\lambda)$ une
$(B_\lambda/\mathfrak J_\lambda)$-algèbre de présentation finie. La conclusion
résulte alors de (11.2.6).
\end{proof}

